% ===================================================================
%  جدول نمبر 1 — اسکول، ثانوی اسکول، جماعت۔
%
%  Ceci n'est PAS une transcription : c'est une traduction, faite en
%  2026, en ourdou d'aujourd'hui. D'ou trois differences avec
%  texto/io et texto/fr, et trois seulement :
%
%   * pas de \nl ni de \cc — la traduction ne suit aucune ligne
%     imprimee, et les lignes de ce fichier ne sont que du confort ;
%   * pas de \begin{VUpage} — elle n'a ni feuillet ni folio ;
%   * le gras se pose sur le TERME seul, ponctuation dehors.
%
%  Les cles « %%K » sont celles de texto/io, macro pour macro pour
%  l'apparat, et les renvois \textsuperscript{(n)} visent les MEMES
%  objets de la planche, dans le MEME ordre.
%
%  LA PREMIERE COLONNE DU RELEVE A DEUX VOISINES DE NATURES
%  DIFFERENTES, ET C'EST TOUTE SON IDENTITE. Les colonnes defendues
%  jusqu'ici l'etaient sur UN front — le cantonais contre le
%  mandarin, le marathi contre le hindi, le coreen contre le hanja,
%  le tamoul contre son propre parle. L'ourdou en a deux :
%
%   * il partage l'ECRITURE avec le pendjabi chahmoukhi (texto/pnb),
%     qui est une autre langue. Meme alphabet, meme nastaliq : la
%     voisine passerait inapercue.
%   * il partage la LANGUE avec le hindi (texto/hi) — meme
%     grammaire, meme lexique de base, au point que les deux se
%     parlent sans traduction. Ils different par l'ecriture, qui
%     empeche toute contamination visible, et par le REGISTRE
%     SAVANT : sanskrit d'un cote, persan et arabe de l'autre.
%
%  kolonoj.py TIENT LES DEUX FRONTS PAR DEUX MOYENS DIFFERENTS. Le
%  premier par la GRAMMAIRE : le genitif pendjabi « دا » n'existe pas
%  en ourdou, qui dit « کا ». Le second par le LEXIQUE : les tatsama
%  du hindi — ودیالیہ, پستک, ادھیاپک — que l'ourdou imprime n'emploie
%  jamais, non parce qu'ils seraient trop savants mais parce que sa
%  colonne savante est l'autre.
%
%  ET DEUX MOTS N'ONT PAS ETE RELEVES, EXPRES. « دی » et « دے » sont
%  les deux autres formes du genitif pendjabi, et ce sont aussi des
%  mots ourdous ordinaires — le passe et l'imperatif de « donner ».
%  Un controle qui les prendrait crierait a chaque page. De meme
%  « اوہ » : pronom pendjabi et interjection ourdoue, une seule
%  graphie pour deux mots. On garde « ایہہ », qui n'est que
%  pendjabi. Un controle qui crie finit desarme.
%
%  L'APPARAT SE SEPARE DE SON VOISIN D'ALPHABET SUR TROIS MOTS.
%  L'ourdou dit جدول pour un tableau la ou le chahmoukhi dit نقشہ.
%  Pour la serie et pour la scene, en revanche, les deux disent le
%  MEME mot arabe — سلسلہ, منظر — et ne se separent que sur
%  l'ordinal qui le compte : دوسرا, تیسرا contre دوجا, تیجا. C'est
%  la premiere fois dans ce releve que deux colonnes partagent le nom
%  d'un role et se distinguent par le nombre.
%
%  LA REGLE MODIFICATEUR–TETE VAUT ICI SANS RETOUCHE. L'ourdou est a
%  tete finale, son verbe finit la phrase, ses postpositions suivent
%  le nom. parigi.py sert donc comme pour le telougou, le coreen et
%  le tamoul ; cinquieme famille, et il n'a toujours pas fallu y
%  toucher.
%
%  LE VOCABULAIRE DE LA CLASSE EST PERSAN ET ARABE JUSQU'AU BOUT :
%
%    تختۂ سیاہ (1)  le tableau noir    جھاڑن (2)   le torchon
%    چاک (3)        la craie           سلیٹ (42)   l'ardoise
%    دوات (25)      l'encrier          جاذب کاغذ (41) le buvard
%    لغت (30)       le dictionnaire    کاپی (28)   le cahier
%    قلمدان (31)    le plumier         پرکار (26)  le compas
%    پیمانہ (50)    la regle           چاقو (56)   le canif
%    بیاض (54)      le carnet          بستہ (57)   le cartable
%    کھونٹی (58)    la patere          انگیٹھی (46) le poele
%    ریحل (65)      le pupitre         صحن (75)    la cour
%    روشندان (78)   le vasistas
%
%  « روشندان » EST EXACTEMENT LE VASISTAS, et il arrive au bon
%  moment. La colonne tamoule venait de casser une serie de
%  trente-huit en devant l'ecrire en deux mots ; l'ourdou le nomme
%  d'un seul — روشن, la lumiere, et دان, ce qui la contient : la
%  petite fenetre haute par ou le jour entre. Trente-neuf fois sur
%  quarante.
%
%  « چپراسی » (77) EST LE CAS LE PLUS INSTRUCTIF DU TABLEAU. L'homme
%  de service de l'ecole s'est dit فرّاش en egyptien, शिपाई en
%  marathe, சிப்பந்தி en tamoul — ces deux derniers du persan
%  sipāhī —, బంట్రోతు en telougou, 수위 en coreen. L'ourdou POSSEDE
%  سپاہی, le mot persan d'origine, et ne l'emploie pas ici : il dit
%  چپراسی, d'apres چپراس, la plaque de cuivre que l'homme porte au
%  ceinturon. La langue qui a prete le mot aux autres en a garde un
%  troisieme pour elle.
%
%  LES QUATRE OPERATIONS ET LA GEOMETRIE SONT DE L'ARABE, ET C'EST LE
%  MEME ARABE QUE L'EUROPE A EMPRUNTE : جمع, تفریق, ضرب, تقسیم ;
%  دائرہ, مربع, مثلث, خط ; حساب, ہندسہ. Sur une planche ou l'ecole est
%  europeenne de bout en bout, le vocabulaire des mathematiques est
%  celui d'ou l'Europe tient les siennes.
%
%  LA GAMME RESTE OCCIDENTALE, comme au telougou, au coreen et au
%  tamoul. L'ourdou a سرگم et ses sept سُر, et il serait tentant de
%  les ecrire ; mais le tableau nomme do-re-mi et pose l'equivalence
%  « = C. D. E. F. G. A. B. ». C'est la gamme d'Europe qui est
%  l'objet.
%
%  TROIS RETOURNEMENTS SUR LES HUIT PAIRES DE parigi.py, ET CE SONT
%  ENCORE LES MEMES TROIS : « stulo (8) en sua katedro (6) »,
%  « nigra tabelo (1) kun la vishilo (2) », « furnelo (46) kun lua
%  tubo (45) ». Cinquieme langue a tete finale, meme partage.
% ===================================================================

%%K t01-apar-0 apar
\VUsaut{2.0mm}
\VUcentre{12.6pt}{120}{وضاحتی کتابچہ}
\VUsaut{3.2mm}
\VUcentre{8.4pt}{160}{منسلک}
\VUsaut{2.6mm}
\VUtitre{15.6pt}{0}{دلماس کے مددگار جدولوں کا}
\VUsaut{3.4mm}
\VUornamento{9pt}
\VUsaut{5.0mm}
\VUcentre{13.2pt}{90}{پہلا سلسلہ}
\VUsaut{3.0mm}
\VUfilet{16mm}
\VUsaut{5.6mm}

%%K t01-apar-1 apar
\VUtitre{15.0pt}{60}{جدول نمبر 1}
\VUsaut{1.4mm}
\VUtitre{11.4pt}{0}{اسکول۔ --- ثانوی اسکول۔ --- جماعت۔}
\VUsaut{2.2mm}
\VUfilet{20mm}
\VUsaut{6.4mm}

%%K t01-tit-1 sub
\VUpk{10.2pt}{40}{عام وضاحت۔}
\VUsaut{3.0mm}

%%K t01-01-1 p
1. --- یہ جدول ایک \VUgras{ثانوی اسکول} کی \VUgras{جماعت} دکھاتا ہے۔ اس
جماعت میں آٹھ \VUgras{میزیں} \textsuperscript{(18)} اور آٹھ
\VUgras{بینچیں} \textsuperscript{(32)} ہیں۔ ہر بینچ پر تین طالب علم
بیٹھتے ہیں۔ \VUgras{استاد} \textsuperscript{(7)} ایک \VUgras{کرسی} پر
\textsuperscript{(8)} بیٹھے ہیں ؛ وہ کرسی اُن کے \VUgras{چبوترے}
\textsuperscript{(6)} پر ہے۔

%%K t01-02-1 p
2. --- چبوترے کی بائیں طرف \VUgras{شیشے کی الماری} \textsuperscript{(15)}
ہے۔ دائیں طرف \VUgras{تختۂ سیاہ} \textsuperscript{(1)} ہے ؛ اُس کے ساتھ
\VUgras{جھاڑن} \textsuperscript{(2)} اور \VUgras{چاک} \textsuperscript{(3)}
رکھے ہیں۔

%%K t01-02-2 p
چبوترے کے اوپر \VUgras{دیواری تختیاں} \textsuperscript{(9, 11,}
\textsuperscript{12)} اور ایک \VUgras{نقشہ} \textsuperscript{(10)}
لٹکے ہوئے ہیں۔

%%K t01-03-1 p
3. --- سب طالب علم اپنی \VUgras{جگہ} \textsuperscript{(17)} پر نہیں
بیٹھے۔ یوآنیس \textsuperscript{(4)} تختۂ سیاہ کے پاس کھڑا ہے ؛ وہ
\VUgras{جھاڑن} پکڑے ہوئے \VUgras{ہندسی شکلیں} مٹا رہا ہے۔

%%K t01-03-2 p
پیٹرس چبوترے کے قریب ہے۔ وہ \VUgras{تحریری کام} \textsuperscript{(5)}
جمع کر کے استاد کے پاس لے جا رہا ہے۔

%%K t01-04-1 p
4. --- \VUgras{یہ} استاد \VUgras{زندہ زبانوں} کے ہیں۔ وہ طالب علموں کو
غیر ملکی زبانیں بولنا، پڑھنا اور لکھنا سکھاتے ہیں \VUgras{(جرمن،
انگریزی، ہسپانوی، اطالوی، روسی} یا \VUgras{فرانسیسی)}۔

%%K t01-05-1 p
5. --- جماعت بہت کشادہ اور بہت روشن ہے۔ پیچھے کی طرف دو
\VUgras{روشندانوں} \textsuperscript{(78)} والی ایک چوڑی \VUgras{کھڑکی}
اور ایک \VUgras{شیشے کا دروازہ} ہے ؛ ان سے ہوا اور روشنی آتی ہے۔

%%K t01-05-2 p
یہ جماعت ٹھنڈی نہیں ہے۔ بیچ میں کمرہ گرم کرنے کے لیے ایک بہت اچھی
\VUgras{انگیٹھی} \textsuperscript{(46)} ہے ؛ اُس کے ساتھ ایک
\VUgras{دھوئیں کی نلی} \textsuperscript{(45)} لگی ہے۔

%%K t01-05-3 p
دیوار میں \VUgras{کھونٹیاں} \textsuperscript{(58)} جڑی ہیں ؛ اُن پر طالب
علموں کی ٹوپیاں اور اوور کوٹ لٹکے ہیں۔

%%K t01-tit-2 sub
\VUsaut{5.2mm}
\VUpk{10.2pt}{40}{کھیل کا میدان}
\VUsaut{3.6mm}

%%K t01-06-1 p
6. --- \VUgras{گھڑی} \textsuperscript{(63)} نو بج کر پانچ منٹ بتا رہی
ہے۔

%%K t01-06-2 p
\VUgras{تفریح کا وقت} \textsuperscript{(76)} ابھی ختم ہوا ہے اور سبق
دوبارہ شروع ہو رہا ہے۔ کھیلنے کی جگہ \VUgras{صحن} \textsuperscript{(75)}
میں ہے۔ کچھ طالب علم کھیلتے ہوئے نظر آتے ہیں۔ \VUgras{نگران استاد}
\textsuperscript{(74)} اُن کی نگرانی کر رہے ہیں۔

%%K t01-07-1 p
7. --- صحن میں اور بھی کئی لوگ نظر آتے ہیں، یعنی \VUgras{معائنہ کار}
\textsuperscript{(73)}، اسکول کے سربراہ \VUgras{(صدر مدرس)}
\textsuperscript{(71)} اور \VUgras{ناظم} \textsuperscript{(72)}۔

%%K t01-07-2 p
معائنہ کار اسکولوں کا معائنہ کرتے ہیں، صدر مدرس ثانوی اسکول چلاتے
ہیں، ناظم پڑھائی کی نگرانی کرتے ہیں۔

%%K t01-07-3 p
چوتھا آدمی \VUgras{چپراسی} \textsuperscript{(77)} ہے۔ وہ غیر حاضری
لکھنے کے لیے بغل میں ایک رجسٹر دبائے پھرتا ہے۔

%%K t01-08-1 p
8. --- صحن کے پیچھے \VUgras{ورزشی آلات} \textsuperscript{(70)} اور
\VUgras{دستکاری} \textsuperscript{(69)} کی کارگاہ ہے۔

%%K t01-08-2 p
جدول کی دائیں طرف \VUgras{موسیقی} اور \VUgras{خاکہ کشی} کی
\VUgras{جماعتیں} نظر آنے لگتی ہیں۔

%%K t01-noto-1 noto
\VUnotes{143.2mm}{%
(*) \textit{بپتسمے کے ناموں} کو بھی لاطینی صورت کے مطابق لکھنے کا
مشورہ دیا جاتا ہے۔ بے شمار اختلافات سے بچنے کا یہی ایک طریقہ ہے۔
پھر \textit{Giovanni, Jean, Johann, Jan, Hans, John, Ivan} وغیرہ میں
سے کس کو چنیں؟ کیا ہم صرف بپتسمے کے ناموں کے لیے الگ لغت بنائیں؟
لاطینی سے اُن کی فاعلی صورت لے کر یوں کہیں :
\VUgras{Ioannes, Ioanna; Iulius, Iulia; Iakobus; Andreas; Lukas.}
(Gram. Kompleta).%
}

%%K t01-tit-3 sub
\VUpk{10.2pt}{40}{تفصیلی وضاحت۔}
\VUsaut{2.4mm}

%%K t01-tit-4 sub
\VUpk{10.2pt}{40}{اچھے طالب علم۔}
\VUsaut{3.0mm}

%%K t01-09-1 p
9. --- چبوترے کی دائیں طرف پہلی بینچ کے طالب علم \VUgras{ہنریکس} (*)،
\VUgras{سٹیفانس} اور \VUgras{کارولس} ہیں۔

%%K t01-09-2 p
ہنریکس بہت توجہ دینے والا اور محنتی ہے ؛ وہ استاد کی بات سنتا ہے۔

اپنی میز پر اُس کے پاس \VUgras{کاپی} \textsuperscript{(28)}،
\VUgras{کتاب} \textsuperscript{(29)}، \VUgras{لغت} \textsuperscript{(30)}
اور \VUgras{قلمدان} \textsuperscript{(31)} ہیں۔ اُس کی کتاب میں بہت سے
\VUgras{صفحے} ہیں، ہر باب میں کئی \VUgras{پیرے} ہیں، اور ہر
\VUgras{سطر} میں بہت سے \VUgras{الفاظ}۔

%%K t01-09-4 p
اُس کے ساتھ والا سٹیفانس \textsuperscript{(24)} پُر جوش ہے۔ وہ اگلی
بار کا سبق اپنی کاپی میں لکھ لیتا ہے۔ اُس کی \VUgras{لکھائی} خوبصورت
ہے۔ اُس کی کتاب بند ہے ؛ ہمیں صرف \VUgras{جلد} \textsuperscript{(27)}
نظر آتی ہے۔ اُس کے پاس اُس کا \VUgras{پرکار} \textsuperscript{(26)}
رکھا ہے۔

%%K t01-09-5 p
کارولس \textsuperscript{(23)} اپنی کتاب ہاتھ سے پکڑے ہوئے ہے ؛ اپنا
سبق دوبارہ پڑھنے کے لیے اُسے کھولتا ہے۔ ہر طالب علم کی طرح اُس کے
سامنے بھی \VUgras{دوات} \textsuperscript{(25)} ہے۔ وہ فرماں بردار ہے۔

%%K t01-10-1 p
10. --- ساتھ والی میز پر \VUgras{لودوویکس}، \VUgras{انتونیس} اور
\VUgras{پالس} ہیں۔ لودوویکس اپنا \VUgras{سبق} \textsuperscript{(22)}
سناتا ہے اور استاد کے \VUgras{سوالوں} کا جواب دیتا ہے۔ انتونیس
\textsuperscript{(21)} بہت بے توجہ ہے ؛ کبھی کبھی استاد اُسے سزا بھی
دیتے ہیں۔ پالس \textsuperscript{(20)} اچھا \VUgras{ساتھی} ہے، ملنسار
اور مددگار۔ وہ بہت توجہ دینے والا ہے۔

%%K t01-tit-5 sub
\VUsaut{4.2mm}
\VUpk{10.2pt}{40}{برے طالب علم۔}
\VUsaut{3.2mm}

%%K t01-11-1 p
11. --- ہنریکس کے پیچھے \VUgras{جارجیس} \textsuperscript{(38)} ہے۔ وہ
برا لڑکا نہیں ہے، مگر \VUgras{باتونی}، بے توجہ اور شرارتی ہے۔ اُس کی
\VUgras{ہلکی مزاجی} اور \VUgras{نافرمانی} سبق کے دوران \VUgras{ہنگامہ}
پیدا کرتی ہیں ؛ اسی لیے اُسے اکثر سخت \VUgras{تنبیہ} ملتی ہے۔ اُس کی
\VUgras{مسودے کی کاپی} \textsuperscript{(35)} میلی ہے ؛ وہ اپنے
\VUgras{قلم} \textsuperscript{(34)} سے اُس پر \VUgras{سیاہی گراتا} ہے،
اسی لیے ہمیشہ اپنی \VUgras{ربڑ} \textsuperscript{(36)} استعمال کرتا
ہے ؛ اُس کا \VUgras{کاپیوں کا بستہ} \textsuperscript{(33)} پھٹا ہوا
ہے، کیونکہ وہ بہت بے پروا ہے۔ پھر وہ زندہ زبانوں کی جماعت میں اپنا
\VUgras{پنسل دان} \textsuperscript{(37)} کیوں لے آتا ہے؟

%%K t01-11-2 p
اُس کے ساتھ والا ایڈمنڈس \textsuperscript{(39)} اُسے پسند نہیں کرتا۔ وہ
سچ مچ کچھ سست ہے، مگر کم گو اور خاموش ہے۔ وہ باتیں نہیں کرتا ؛ البتہ
سننے کے بجائے اپنی \VUgras{قرأت کی کتاب} کی \VUgras{کہانیاں} پڑھنا
زیادہ پسند کرتا ہے۔

%%K t01-11-3 p
اس بینچ کا تیسرا طالب علم \VUgras{لودوویکس} \textsuperscript{(40)} ہے۔
وہ محنتی ہے۔ سفید \VUgras{کاغذ} پر لکھتا ہے۔ اپنے سامنے اپنا
\VUgras{جاذب کاغذ} \textsuperscript{(41)} رکھ لیا ہے۔

%%K t01-12-1 p
12. --- استاد کی بائیں طرف دوسری بینچ کے طالب علم بہت چھوٹے ہیں۔
پہلا، ننھا \VUgras{ایمیلیس}، ابھی اچھی طرح لکھنا نہیں جانتا ؛ وہ اپنی
\VUgras{سلیٹ} \textsuperscript{(42)}، اپنی \VUgras{سلیٹی} اور اپنا
\VUgras{اسفنج} \textsuperscript{(43)} ساتھ لاتا ہے۔

%%K t01-12-2 p
اُس کے پاس \VUgras{اگستس} اور \VUgras{برنارڈس} \textsuperscript{(44)}
ہیں۔ یہ اچھی طرح کام کرتا ہے، مگر اُس کا \VUgras{لباس} بہت بے ترتیب
ہے۔

%%K t01-13-1 p
13. --- اُن کے پیچھے تین \VUgras{سست لڑکے} ہیں۔ \VUgras{البرٹس} اپنے
سبق یاد نہیں کرتا۔ \VUgras{یاکوبس} \textsuperscript{(47)} سننے کے بجائے
سوتا ہے۔ اُس کی \VUgras{سستی} اُسے \VUgras{سرزنش} اور \VUgras{سزا}
دلاتی ہے۔ آخر میں \VUgras{کرسٹیانس} \textsuperscript{(48)} بہت غیر
سنجیدہ ہے۔ وہ برا \VUgras{برتاؤ} کرتا ہے اور سدھرنے کی ذرا کوشش نہیں
کرتا۔

%%K t01-14-1 p
14. --- اُس کے برعکس اُس کے ساتھی خوش اطوار ہیں۔ \VUgras{ارنسٹس}
\textsuperscript{(51)} ترتیب اور سلیقے کو پسند کرتا ہے ؛ ہمیشہ اپنا
\VUgras{پیمانہ} \textsuperscript{(50)} اور اپنی \VUgras{پنسل}
\textsuperscript{(49)} استعمال کرتا ہے۔ وہ \VUgras{روزانہ آنے والا
طالب علم} ہے۔

%%K t01-14-2 p
\VUgras{فرانسسکس} \textsuperscript{(52)} \VUgras{اقامتی طالب علم} ہے ؛
وردی پہنتا ہے اور ثانوی اسکول ہی میں کھاتا اور سوتا ہے۔ وہ اچھا
\VUgras{ساتھی} ہے۔

%%K t01-14-3 p
مورسیس اپنے \VUgras{چاقو} \textsuperscript{(56)} سے اپنی \VUgras{پنسل}
چھیلتا ہے۔ وہ بہت محتاط ہے۔

%%K t01-14-4 p
وہ اپنی سب کتابیں ساتھ لاتا ہے — اپنی \VUgras{قواعد کی کتاب}
\textsuperscript{(55)}، اپنی \VUgras{بیاض} \textsuperscript{(54)}، حتیٰ
کہ اپنا \VUgras{تختۂ تحریر} \textsuperscript{(53)} بھی۔ اُس کا
\VUgras{بستہ} \textsuperscript{(57)} اُس کے پیچھے فرش پر رکھا ہے۔

%%K t01-14-5 p
جماعت کے پیچھے کی دونوں میزیں \VUgras{اچھے طالب علم}
\textsuperscript{(19)} استعمال کرتے ہیں۔

%%K t01-tit-6 sub
\VUsaut{4.6mm}
\VUpk{10.2pt}{40}{خاکہ کشی اور موسیقی۔}
\VUsaut{3.4mm}

%%K t01-15-1 p
15. --- \VUgras{خاکہ کشی کے کمرے} میں دیوار پر لٹکے خوبصورت
\VUgras{نمونے} ہیں، اور چھت سے \VUgras{سایہ بان} والا ایک
\VUgras{لیمپ} \textsuperscript{(62)} لٹکا ہے۔ \VUgras{استاد}
\textsuperscript{(60)} بہت صابر ہیں۔ وہ طالب علموں کے \VUgras{خاکے}
\textsuperscript{(61)} درست کرتے ہیں اور اُنھیں اصل دیکھ کر
\VUgras{نیم تن مجسمہ} \textsuperscript{(59)} بنانا سکھاتے ہیں۔

%%K t01-16-1 p
16. --- \VUgras{موسیقی کا کمرہ} بہت دلچسپ لگتا ہے۔ وہاں
\VUgras{موسیقی} پڑھنا، \VUgras{سرگم کی لکیروں} پر لکھے
\VUgras{سُر} \textsuperscript{(67)} گانا (دو، رے، می، فا، سول، لا،
سی\,=\,C. D. E. F. G. A. B.)، \VUgras{کلیدیں} پہچاننا اور خوبصورت
\VUgras{دھنیں} گانا سکھایا جاتا ہے۔ موسیقی کے کاغذ \VUgras{ریحل}
\textsuperscript{(65)} پر رکھے جاتے ہیں۔ ایک طالب علم \VUgras{پیانو
بجاتا ہے} \textsuperscript{(64)} اور \VUgras{موسیقی کے استاد}
\textsuperscript{(66)} \VUgras{تال دیتے ہیں} \textsuperscript{(68)}۔

%%K t01-17-1 p
17. --- \VUgras{دستکاری} \textsuperscript{(69)} کی \VUgras{کارگاہ} کا
ایک کونہ نظر آنے لگتا ہے ؛ وہاں لڑکے لکڑی رندنا اور لوہا ریتنا سیکھ
سکتے ہیں۔ یہ ہر ثانوی اسکول میں نہیں ہوتا۔

%%K t01-tit-7 sub
\VUsaut{5.0mm}
\VUpk{10.2pt}{40}{سائنس کی جماعت۔}
\VUsaut{3.6mm}

%%K t01-18-1 p
18. --- ریاضی کے استاد طالب علموں کو \VUgras{ہندسہ}، \VUgras{حساب}،
\VUgras{شمار} اور \VUgras{میٹری نظام} پڑھاتے ہیں۔ تختۂ سیاہ پر ہمیں
اہم ہندسی شکلیں نظر آتی ہیں : \VUgras{دائرہ} \textsuperscript{(a)}،
\VUgras{مربع} \textsuperscript{(b)}، \VUgras{مثلث} \textsuperscript{(c)}،
\VUgras{خط} \textsuperscript{(d)}۔

%%K t01-18-2 p
ہمیں ایک \VUgras{سوال} بھی نظر آتا ہے۔ وہ نہ \VUgras{جمع} ہے، نہ
\VUgras{تفریق}، نہ \VUgras{ضرب} ؛ وہ تو \VUgras{تقسیم}
\textsuperscript{(e)} ہے۔

%%K t01-19-1 p
19. --- میٹری نظام کے جدول پر ہمیں یہ نظر آتے ہیں : \VUgras{میٹر}
\textsuperscript{(a)}، \VUgras{ترازو} \textsuperscript{(b)} اور اُس کے
\VUgras{باٹ}، \VUgras{لیٹر} \textsuperscript{(e)}، \VUgras{ڈیکا لیٹر}
\textsuperscript{(f)}، \VUgras{ڈیسی لیٹر}، \VUgras{مکعب میٹر}
\textsuperscript{(d)}۔

%%K t01-19-2 p
طالب علم \VUgras{علومِ فطرت} \textsuperscript{(11)}، علمِ حیوانات
\textsuperscript{(a)}، علمِ نباتات \textsuperscript{(b)}، علمِ ارضیات
\textsuperscript{(c)} اور \VUgras{تاریخ} \textsuperscript{(12)} بھی
پڑھتے ہیں۔

%%K t01-19-3 p
الماری میں \VUgras{طبیعیات} \textsuperscript{(15)} اور
\VUgras{کیمیا} \textsuperscript{(16)} کے آلات ہیں۔

%%K t01-19-4 p
الماری کے اوپر \VUgras{کرہ} \textsuperscript{(14)}، \VUgras{مخروط} اور
\VUgras{کرۂ ارض} \textsuperscript{(13)} رکھے ہیں۔

%%K t01-19-5 p
جدولوں کے اوپر ایک \VUgras{نقشہ} لٹکا ہے جو دنیا کے پانچ حصے دکھاتا
ہے : \VUgras{یورپ} \textsuperscript{(g)}، \VUgras{ایشیا}
\textsuperscript{(e)}، \VUgras{افریقہ} \textsuperscript{(f)}،
\VUgras{امریکہ} \textsuperscript{(ab)}، \VUgras{اوقیانوسیہ}
\textsuperscript{(h)}، اور دو بڑے سمندر — \VUgras{بحرالکاہل}
\textsuperscript{(c)} اور \VUgras{بحرِ اوقیانوس} \textsuperscript{(d)}۔

\VUsaut{6.0mm}
\VUfilet{22mm}
